\chapter{Implementations}

\section{Choose the model first}
A state must determine every future legal move and cost. A cell alone is
insufficient when keys, parity, orientation, or remaining wall breaks matter.
Use $(r,c,\mathit{mask})$, $(r,c,t\bmod p)$, or the required full state.
Count states times transitions before allocating memory.

\textbf{Shortest paths.} Unit costs: BFS. Many sources: one BFS with every
source at distance zero. Costs in $\{0,1\}$: deque-based 0--1 BFS.
Nonnegative costs: Dijkstra. DAG: topological DP, even with negative costs.
Negative edges in a general graph: use Bellman--Ford, not Dijkstra.

\textbf{Connectivity.} Static undirected graph: components. Adding edges:
DSU. Cheapest connected edge set: Kruskal. Directed mutual reachability:
SCC. DSU gives neither paths nor shortest distances.

\textbf{Ranges.} Immutable sums: prefix sums. Offline additions: differences.
Point additions/range sums: Fenwick. Point assignment/associative range
aggregate: segment tree. Range updates require a compatible lazy structure.

\textbf{Sequences.} Nonnegative contiguous sums with monotone feasibility:
two pointers. First feasible numeric answer: binary search after proving
monotonicity. Repeated histories with the same future options: DP.

\subsection{Conventions and numeric limits}
All companion APIs use zero-based vertices and half-open ranges $[l,r)$.
Convert one-based inclusive $[L,R]$ to $[L-1,R)$. Empty range: $l=r$.
Prefix arrays have $n+1$ entries. Weighted unreachable distance is
\verb@INF=LLONG_MAX/4@; all represented finite distances must be smaller.
Other sums/products still need their own overflow bound; cast before
multiplication. Graph endpoints and sources must be valid. Grid rows must
have equal lengths; $hw$ must fit \texttt{int}. Reset storage per test case.
Module snippets share the \texttt{finals\_core} namespace and the types
from the common header. Copy their dependencies too.

\kactlimport{Prelude.h}

\section{Unweighted graphs}
\subsection{BFS, parents, and multiple sources}
\verb@bfs(g,{s})@ returns distances in edges, unreachable $-1$, and parents.
Mark a vertex when enqueuing it. The first discovery is shortest because
FIFO order processes distance $d$ before $d+1$. With several sources,
parent $-1$ identifies whichever source supplies the returned shortest path.
Duplicate sources are ignored; no sources leaves every vertex unreachable.
Time $O(n+m+|S|)$, extra memory $O(n)$.

\verb@restore_path(result,t)@ returns source through $t$, or an empty vector
if unreachable. Source equal to target gives one vertex and zero edges.
Adjacency order chooses ties; lexicographic minimality is not automatic.
\begin{lstlisting}
auto a = bfs(g, {source});
auto p = restore_path(a, target);
if (p.empty()) { /* impossible */ }
else { /* p.size()-1 edges */ }
\end{lstlisting}
\kactlimport{BFS.h}

\subsection{Connected components}
\verb@components(g)@ is iterative and assumes an undirected graph.
It returns component IDs, representatives, sizes, and the count.
Isolated vertices count. Connect representative $0$ to every other
representative to add the minimum $c-1$ roads. The lower bound follows
because one edge decreases the component count by at most one.
Time $O(n+m)$, extra memory $O(n)$.
For directed input this routine computes neither SCCs nor weak components;
use SCC or explicitly ignore edge directions, respectively.
\kactlimport{Components.h}

\subsection{Grid BFS and time-dependent movement}
\verb@grid_bfs(board,{{sr,sc}},'#')@ uses four-neighbor moves and treats
all non-wall characters as passable. ID of $(r,c)$ is $rw+c$;
recover $(r,c)$ as \verb@id/w,id%w@. Parents use these same IDs.
For successive path cells, coordinate differences identify U/D/L/R.
Time and memory $O(hw)$; no explicit adjacency lists are needed.
A wall source is ignored. A source outside the grid is invalid.

For spreading hazards, compute all hazard arrival times first with a
multi-source BFS. In a \emph{separate constrained} player BFS, enter $v$
at time $t$ only when legal, often $t<\mathit{hazard}[v]$.
The helper below does not enforce hazard constraints. Simultaneous arrival
may be forbidden. Recheck boundary escape and start-cell rules.
\kactlimport{GridBFS.h}

\section{Weighted paths and dependencies}
\subsection{0--1 BFS}
Use adjacency entries \verb@{neighbor,weight}@. A strict improvement with
weight zero enters the deque front; weight one enters the back. A vertex
may improve after its first discovery: a permanent visited flag is wrong.
The returned parent chain is valid after checking distance $\ne\texttt{INF}$.
Time $O(n+m)$, space $O(n+m)$ including queued entries.
Example reduction: an allowed directed edge costs zero; traversing its
reverse costs one. Use this only if the statement permits that move.
% WeightedPaths.h follows the Dijkstra contract below.

\subsection{Dijkstra}
\verb@dijkstra(g,s)@ requires nonnegative weights and returns distances and
parents. Ignore stale heap entries. The implementation skips candidate
distances at least \texttt{INF} before adding, preventing that addition's
overflow but leaving such distances unrepresentable.
Time $O((n+m)\log(n+m+1))$, space $O(n+m)$.
\begin{lstlisting}
auto a = dijkstra(g, source);
if (a.dist[target] == INF) {
  /* unreachable */
} else { /* answer a.dist[target] */ }
\end{lstlisting}
Zero-weight cycles are valid for distances. Counting shortest paths by
adding ways during heap processing needs extra reasoning in their presence.
For weighted path reconstruction, follow \texttt{parent} to $-1$ and reverse;
\texttt{restore\_path} accepts the unweighted \texttt{Search} type only.
\kactlimport{WeightedPaths.h}

\subsection{Topological order and DAG DP}
Kahn's algorithm repeatedly removes an indegree-zero vertex.
\verb@topological_order(g)@ returns a full order, or empty on a cycle.
Check \verb@order.size()==g.size()@, including the valid empty graph.
Parallel directed edges contribute separately to indegrees.
Time $O(n+m)$, space $O(n)$.

Longest path from $s$: initialize $dp[s]=0$, others $-\infty$.
In topological order, for reachable $u$, apply
$dp[v]\gets\max(dp[v],dp[u]+w_{uv})$ and save its parent on improvement.
Unreachable is different from length zero. For shortest DAG paths use
minimum; for path counts use addition modulo $M$. Parallel edges count
as separate paths unless deduplicated. Never apply DAG DP after rejecting
a cyclic graph. An ordinary vertex index order is not a topological order.
\kactlimport{TopoSort.h}

\section{Connectivity and trees}
\subsection{DSU and Kruskal}
\verb@DSU(n)@ starts with singleton sets. \verb@unite(u,v)@ returns whether
it merged sets; \verb@find(u)@ gives a representative and \verb@size(u)@
the current set size. Representatives may change. Amortized time per
operation $O(\alpha(n))$, memory $O(n)$. Arbitrary deletions need a different
method; some offline deletion tasks become additions when reversed.

\verb@kruskal(n,edges)@ accepts undirected \verb@Edge{u,v,w}@, stored once
per edge. It returns selected edges, total weight, and remaining component
count. Negative weights, parallel edges, and loops are allowed; loops
never merge sets. Connected nonempty input has one component and $n-1$
selected edges. Otherwise the answer is a minimum spanning \emph{forest}.
Time $O(m\log m)$; prove that total weight fits \texttt{ll}.
MST minimizes selected edge sum, not distances from a source.
\kactlimport{DSUKruskal.h}

\subsection{Iterative tree DP}
\verb@root_tree(tree,root)@ assumes a nonempty connected undirected tree.
It returns parents, depths, parent-before-child order, and subtree sizes.
Root is its own parent; size includes the vertex itself. Reverse order is
suitable for child-to-parent DP. Time and memory $O(n)$, no recursion stack.
On cyclic input this merely gives a spanning tree: the caller must verify
the tree assumption. Its order is not an Euler tour; subtree vertices need
not form an interval. For range queries compute entry/exit times separately.

Maximum independent set state: $T[u]$ selects $u$, $S[u]$ skips it.
Initialize $T[u]=1$, $S[u]=0$. For each child $v$,
\[
T[u]\mathrel{+}=S[v],\qquad
S[u]\mathrel{+}=\max(T[v],S[v]).
\]
Answer $\max(T[root],S[root])$. Child subtrees are independent because a tree
has no cross edges. For weighted selection initialize $T[u]$ to its weight.
Save the root decision; when reconstructing, selecting a parent forces every
child skipped, otherwise choose each child's better state.

\subsection{Binary lifting LCA}
\verb@LCA a(info)@ uses the rooted tree data. \verb@a.jump(u,k)@ returns
the $k$th ancestor or $-1$ above the root. \verb@a.query(u,v)@ is the lowest
common ancestor. \verb@a.distance(u,v)@ counts edges, not weights.
Preprocess $O(n\log n)$ time/memory, queries $O(\log n)$.
For weighted trees, accumulate root distances separately and use
$D[u]+D[v]-2D[\mathrm{lca}(u,v)]$. Cross-component forest queries are invalid.
\kactlimport{TreeLCA.h}

\section{Intervals and monotonicity}
\subsection{Coordinate compression and prefix sums}
\verb@compressed_values(a)@ sorts and removes duplicates.
\verb@rank_of(values,x)@ is the lower-bound insertion index. It represents
$x$ exactly only if $x$ occurs in the compressed list. Compression preserves
order/equality, not distances. Keep original coordinates for lengths/areas.
Build $O(n\log n)$, lookup $O(\log n)$.

\verb@prefix_sums(a)@ gives $P[0]=0$, $P[i+1]=P[i]+a[i]$;
interval sum $[l,r)$ is $P[r]-P[l]$. Negative entries are allowed.
Build $O(n)$, query $O(1)$. Online updates invalidate static prefixes.
\kactlimport{CompressionPrefix.h}

\subsection{Offline difference array}
Allocate \verb@vector<ll> diff(n+1)@.
\verb@difference_add(diff,l,r,x)@ adds $x$ on $[l,r)$.
After all updates, \verb@materialize_difference(diff)@ gives $n$ values.
For an existing array, add these increments to it. Updates cost $O(1)$,
materialization $O(n)$. The extra sentinel is needed when $r=n$.
This method does not support interleaved queries without rebuilding.
\kactlimport{DifferenceArray.h}

\subsection{Sliding window}
\verb@count_subarrays_at_most(a,k)@ requires all $a_i\ge0$ and counts
nonempty contiguous subarrays of sum at most $k$. It returns zero for $k<0$.
For each right endpoint remove left values until feasible; then $r-l+1$
starts are valid. Both pointers move only forward: $O(n)$ time, $O(1)$ space.
Exact sum $k$ equals $\mathrm{atMost}(k)-\mathrm{atMost}(k-1)$.
With negative entries, this monotonicity fails. For arbitrary signed values
and exact sum, count previous prefixes equal to $P-k$; initialize count of
prefix zero to one and query before inserting the current prefix.
\kactlimport{SlidingWindow.h}

\subsection{Binary search on the answer}
\verb@first_true(lo,hi,ok)@ searches nonnegative $[lo,hi)$ with a
false-then-true predicate. It returns original $hi$ when no tested value
works and never evaluates \verb@ok(hi)@. If $hi$ is a proven feasible
sentinel, returning it is also a valid answer. State that convention.
Time $O(\log(hi-lo+1)\cdot\mathrm{cost}(ok))$.

Prove ``if capacity $x$ works, every larger capacity works'' before searching.
A greedy feasibility routine needs a separate proof. Watch overflow inside
\texttt{ok}; safe midpoint arithmetic does not protect products there.
\kactlimport{BinarySearch.h}

\section{Dynamic range structures}
\subsection{Fenwick sums and order statistics}
\verb@Fenwick(n)@ starts with zeros. \verb@add(i,d)@ adds at index $i$;
\verb@prefix(r)@ sums $[0,r)$ and \verb@sum(l,r)@ sums $[l,r)$.
Assignment requires \verb@add(i,newValue-oldValue)@ and updating the saved
old value. Each operation costs $O(\log n)$, memory $O(n)$.
Negative entries work for sums.

\verb@lower_bound(k)@ requires $k>0$ and nonnegative represented entries.
It returns the zero-based position whose inclusive prefix first reaches $k$,
or $n$ if total $<k$. Frequencies $[0,2,0,3]$ give positions $1$ for ranks
1--2, $3$ for ranks 3--5. Remove a selected item with \verb@add(pos,-1)@,
keeping frequencies nonnegative.

Inversions: scan left to right, query
$seen-\mathrm{prefix}(rank+1)$, then add one at the current rank.
Query-before-update excludes the current value; $rank+1$ excludes equals
from the inversion count.
\kactlimport{Fenwick.h}

\subsection{Iterative segment tree}
\verb@SegTree<T,Op>(a,identity)@ supports \verb@set(i,x)@ and
\verb@query(l,r)@ on $[l,r)$. Build/storage $O(n)$, operations $O(\log n)$.
The operation must be associative with a two-sided identity. Order is
preserved, so commutativity is unnecessary. Empty range returns identity.
\begin{lstlisting}
SegTree<ll,Sum> sums(a,0LL);
const ll HI = std::numeric_limits<ll>::max();
SegTree<ll,Minimum> mins(a,HI);
sums.set(i,x);
ll answer = sums.query(l,r);
\end{lstlisting}
Identities: sum $0$, minimum $+\infty$, maximum $-\infty$, GCD $0$.
The two example trees are independent; changing one does not change the other.
A maximum-subarray node needs total, best prefix, best suffix, and best
subarray, with consistent empty-subarray semantics. Scalar sums are insufficient.
For range updates use lazy propagation with proved tag composition;
this code supports point assignment only.
\kactlimport{SegmentTree.h}

\section{DP state and reconstruction}
\subsection{A state before a template}
Write: exact state meaning, transitions, base cases, impossible value,
dependency order, state count times transitions, answer state, witness method.
Merge histories only when their future legal choices and effects are identical.
A DP with huge numeric capacity may still time out; consider value-based DP,
meet-in-the-middle, or a smaller state before allocating it.

\subsection{Worked 0--1 knapsack state}
$D_i[c]$ is the maximum value using the first $i$ items with total weight
\emph{at most} $c$. Base $D_0[c]=0$ and
\[
D_{i+1}[c]=\max(D_i[c],D_i[c-w_i]+v_i)
\quad\text{when }w_i\le c.
\]
Otherwise keep $D_i[c]$. Compress to one row and scan capacities downward,
so the right side still belongs to the previous item layer.
For weights $[2,3,4]$, values $[4,5,7]$, capacity 5, answer 9 uses items
of weight 2 and 3. Upward scanning would incorrectly reuse an item.

\verb@knapsack01(w,v,C)@ permits nonnegative weights, including zero;
\verb@knapsack_unbounded(w,v,C)@ requires strictly positive weights and scans
upward to reuse items. Both maximize value at most capacity and allow taking
nothing; time $O(nC)$, memory $O(C)$. Negative values can be skipped.
A positive-value zero-weight unbounded item makes the optimum infinite.

For \emph{exact} weight, set $D[0]=0$ and others $-\infty$, and skip
unreachable predecessors. Zero initialization everywhere would invent
reachable states. Counting uses addition instead of maximum and needs its
own loop-order argument.

\subsection{Reconstructing selected items}
A mutable parent pointer in overwritten one-row DP can produce an invalid
witness. Simple safe method: retain all rows $D_i[c]$. Starting from
$(i,c)=(n,C)$, if $D_i[c]=D_{i-1}[c]$, skip item $i-1$; otherwise record it
and set $c\gets c-w_{i-1}$. Decrease $i$ in either case.
Zero-weight items remain safe because the item index decreases.
This uses $O(nC)$ memory; if that is too large, use immutable predecessor
records or a separately proved reconstruction scheme.
\kactlimport{Knapsack.h}

\subsection{Counting order and subsequences}
Coins outside, sums ascending: unordered combinations. Sums outside,
last coin inside: ordered sequences. With coins $\{1,2\}$ and total 3,
there are 2 combinations but 3 ordered sequences. Deduplicate denominations
unless equal values intentionally represent distinct types.

\verb@lis_indices(a)@ reconstructs one longest \emph{strictly} increasing
subsequence in $O(n\log n)$ time, $O(n)$ memory. Empty input returns empty.
Tails store minimum achievable last values, not necessarily one actual
subsequence; predecessor indices recover the witness.
For nondecreasing subsequences, use \texttt{upper\_bound} instead of
\texttt{lower\_bound}. Repeated equal values must not extend strict LIS.
\kactlimport{LIS.h}


\subsection{Digit patterns and remainder state}
\verb@finals_dp::complete_digits(pattern,m)@ finds the lexicographically
smallest completion divisible by positive $m$, or \texttt{nullopt}.
The pattern is nonempty and contains digits or \texttt{?}. Leading zero is
forbidden by default, even for one digit; pass \texttt{false} if zero or
fixed-width leading zeros are valid. Time $O(10Lm)$, memory $O(Lm)$.
Check that product before allocation; this is not suitable for huge $m$.

State $can[pos][r]$ means the suffix can finish a number whose already read
prefix has remainder $r$. Appending digit $d$ gives
$r'=(10r+d)\bmod m$. Base $can[L][0]=1$, other final remainders false.
Fill positions backward. Reconstruct forward by trying allowed digits
0 through 9 and taking the first whose next state can finish. This is
lexicographically smallest because every smaller feasible next digit is
tried first, while the suffix table guarantees completion.
For pattern \texttt{?2} and $m=7$, answer \texttt{42}.

Counting integers $\le N$ is a different digit DP: add a \texttt{tight}
state controlling the next allowed digit, and usually a \texttt{started}
state for leading zeros. Do not treat these zeros as real digits when
counting digit occurrences unless the statement requires it.
\kactlimport{DigitDP.h}

\section{Before submitting}
Check edges versus vertices, inclusive versus exclusive endpoints, strict
versus non-strict comparison, impossible states, and numeric bounds.
Try unreachable target, start equal target, zero weight, repeated coordinate,
empty range, last index, multiple test cases, and one item that fits twice.
Tests verify code mechanics; a reduction still needs a proof.
